% AI-produced by GPT-6 Astra (OpenAI, ChatGPT Pro mode).
\subsection{The full component image, with its endpoint condition retained}

\begin{saresult}{A.137}{Proved under stated hypotheses}{A complete component-image criterion with its endpoint and cohomology obstructions retained}
Let $B$ be a closed compact real PL four-manifold. Write
$c_A:A_{\mathrm{PL}}\to A_{\mathrm{TOP}}$ and
$c_P:P_{\mathrm{PL}}\to P_{\mathrm{TOP}}$ for the ordinary-family
comparisons of A.135, and $e$ for top restriction. Define the subgroup
\[
 I_B=(\pi_0e)^{-1}
       \bigl(\operatorname{im}(\pi_0c_A)\bigr)
       \subseteq\pi_0P_{\mathrm{TOP}}(B),
 \qquad
 J_B=\operatorname{im}(\pi_0\beta_B)
       \subseteq H^2(B;\mathbb Z/2),                    \tag{A.137.1}
\]
where $\beta_B$ is the map specified in A.136. There is a
well-defined homomorphism
\[
 \lambda_B:I_B\longrightarrow H^2(B;\mathbb Z/2)/J_B
 \quad\hbox{such that}\quad
 \ker\lambda_B=\operatorname{im}(\pi_0c_P).             \tag{A.137.2}
\]
No surjectivity of $\lambda_B$ is asserted.
In particular $\operatorname{im}(\pi_0c_P)$ is normal in $I_B$ and
\[
 I_B/\operatorname{im}(\pi_0c_P)
     \lhook\joinrel\longrightarrow H^2(B;\mathbb Z/2)/J_B.
                                                               \tag{A.137.2a}
\]
Thus this quotient is a finite elementary abelian $2$-group, of
order at most $2^{\dim_{\mathbb Z/2}H^2(B;\mathbb Z/2)}$.

Here is an explicit definition. Given $h\in P_{\mathrm{TOP}}(B)$
whose component belongs to $I_B$, choose a TOP isotopy $\gamma$
from its top map $f=e(h)$ to a PL map $a$. Lift this endpoint change
to a bottom-fixed isotopy of the cylinder and let $h_\gamma$ be its
final composite with $h$, so that $e(h_\gamma)=a$. Put
\[
 \kappa(h_\gamma)
   =[\,(h_\gamma)_*\alpha_N\,]-[\,\alpha_N\,]
   \in H^3(N,\partial N;\mathbb Z/2)
   \cong H^2(B;\mathbb Z/2),\qquad N=B\times I,         \tag{A.137.3}
\]
using the relative PL-structure classification with the standard
structure $\alpha_N$ as origin. Then
$\lambda_B([h])$ is the class of (A.137.3) modulo $J_B$.
Product collars are imposed before taking the relative class.

Consequently $h$ is TOP-isotopic, through bottom-fixed ordinary
concordances with free top, to a PL concordance if and only if its top
is TOP-isotopic to a PL map and the class (A.137.3) is zero modulo
$J_B$. If $H^2(B;\mathbb Z/2)=0$, the endpoint condition is the only
condition. The empty-base case has trivial groups throughout.

This is a criterion for the complete image on components. It leaves
both the ordinary four-dimensional endpoint-image predicate and
$J_B$ unevaluated in general. It does not claim that vanishing of a
stable boundary obstruction makes a four-dimensional homeomorphism
isotopic to PL, or that every TOP concordance component is represented
by a PL concordance.
\end{saresult}
\begin{proof}
All cohomology in this proof has coefficients $\mathbb Z/2$.
Composition is written so that $FG$ means $F\circ G$.
The common-product-collar normalization of A.136 applies to each
finite family without changing either boundary restriction. We use
it whenever a relative structure on $N$ is mentioned.

\textbf{The relative class when the top is already PL.}
Let $D_B$ be the subgroup of ordinary TOP concordances whose top map
is literally PL. A member $F$ of $D_B$ is PL on both boundary
components; after collar normalization its pushed-forward structure
$F_*\alpha_N$ agrees with $\alpha_N$ on neighbourhoods of both
boundaries. The open five-dimensional relative classification used
in A.135 therefore defines $\kappa(F)$ by (A.137.3). In particular
this use has the \emph{entire} four-dimensional boundary controlled.
The relevant original inputs are Kirby--Siebenmann, Essay~V,
relative (1.5), p.~229, relative classification (3.2), pp.~240--241,
and the chosen-lift description in \S5.6, pp.~252--253
\cite{KirbySiebenmann1977}.

We record the exact zero test:
\[
 \kappa(F)=0
 \quad\Longleftrightarrow\quad
 F\text{ is TOP-isotopic rel. }\partial N
      \text{ to a PL map}.                             \tag{A.137.4}
\]
Indeed zero means that $F_*\alpha_N$ belongs to the component of the
standard relative structure. A path from $\alpha_N$ to this
structure lifts through the relative automorphism orbit map, starting
at the identity, by the orbit-component and Kan-lifting statements
proved from the cited source in A.135. Let its lifted endpoint be
$g$, so that $g_*\alpha_N=F_*\alpha_N$ and $g$ is joined to the
identity through TOP maps fixed near $\partial N$. Then $g^{-1}F$
is PL, and the lifted path gives the required isotopy from $F$ to
$g^{-1}F$. Conversely such an isotopy supplies a path of relative
structures from $F_*\alpha_N$ to the standard structure. This also
shows that $\kappa$ is constant on paths in $D_B$ with their boundary
maps fixed.

\textbf{Additivity and the absence of a coefficient action.}
The difference class of two relative PL lifts is additive and
natural. One can see both statements in the universal
Kirby--Siebenmann square of Essay~V, \S5.6: the chosen standard lift
identifies the relative lift space with
$\operatorname{Map}_*(N/\partial N,K(\mathbb Z/2,3))$;
differences of lifts are added in this mapping space. For a
boundary-PL homeomorphism, naturality gives the affine cocycle rule
\[
 \kappa(FG)=\kappa(F)+F_*\kappa(G),\qquad F,G\in D_B.   \tag{A.137.5}
\]
Here $F_*$ is the action on the relative cohomology group; using
pullbacks instead merely reverses the convention and has no effect
on the conclusion below.

Every bottom-fixed $F$ acts trivially on $H^3(N,\partial N)$.
Its top map $a$ is homotopic to the identity: the homotopy is
$(x,t)\mapsto\operatorname{pr}_B F(x,t)$. The restriction maps
$H^j(N)\to H^j(\partial N)$ are the diagonal maps
$H^j(B)\to H^j(B)\oplus H^j(B)$ and are injective. The exact
cohomology sequence therefore identifies
\[
 H^3(N,\partial N)
  \cong\operatorname{coker}
       \bigl(H^2(B)\xrightarrow{(1,1)}
                      H^2(B)\oplus H^2(B)\bigr).       \tag{A.137.6}
\]
On the two boundary factors $F$ induces respectively the identity
and $a^*=1$. Naturality of the connecting map in (A.137.6) proves
the assertion. Thus (A.137.5) reduces to
\[
                 \kappa(FG)=\kappa(F)+\kappa(G).       \tag{A.137.7}
\]
The argument works componentwise when $B$ is disconnected, since a
bottom-fixed concordance preserves each cylinder component.

\textbf{Traces and the subgroup $J_B$.}
For a TOP path $\delta_t$ from the identity to a PL map $b$, make
the path constant near $t=0,1$ and form its trace
\[
 T_\delta(x,t)=(\delta_t(x),t).                         \tag{A.137.8}
\]
It belongs to $D_B$. It is in the identity component of the ordinary
free-top concordance group, by the explicit isotopy
$(x,t)\mapsto(\delta_{s t}(x),t)$, $0\le s\le1$.
These formulas are TOP formulas; no joint PL property is needed
or asserted for an arbitrary TOP path. By the geometric description
in A.136,
\[
 \kappa(T_\delta)=\pi_0\beta_B([b,\delta]).             \tag{A.137.9}
\]
Every element of $J_B$ arises in this way, since the path model of
the homotopy fibre $F_A$ consists of such pairs $(b,\delta)$.
The image is a subgroup: $\beta_B$ is a loop map under the
identification $F_A\simeq\Omega\mathcal S_B$ in A.136.

\textbf{Changing the top and comparing all choices.}
Let $f=e(h)$ and let $\gamma_t$ be a TOP isotopy from $f$ to a PL
map $a$, constant near its two parameter endpoints. Set
\[
 E_{\gamma,s}(x,t)
       =(\gamma_{s t}f^{-1}(x),t),\qquad
 h_\gamma=E_{\gamma,1}h.                               \tag{A.137.10}
\]
This is a continuous isotopy through bottom-fixed concordances,
$E_{\gamma,0}=1$, and $h_\gamma$ has top $a$. Its inverse is
obtained slice by slice from $f\gamma_{s t}^{-1}$, so no extension
or existence assertion is hidden in this formula.
It is the explicit path-lifting construction for top evaluation
from A.135. At $s=1$ its collars are products; common collar
normalization may be applied to the entire path if desired.
Two lifts of the same endpoint isotopy, both starting at the
identity, differ by a path in $H_{\mathrm{TOP}}$ starting at the
identity. Their final composites with $h$ are therefore isotopic
with both boundary maps fixed and have the same $\kappa$. Thus the
arbitrary choice of lift in the statement is harmless.

For another choice $\eta:f\leadsto b$ with $b$ PL, exact composition
gives
\[
 h_\eta h_\gamma^{-1}(x,t)
           =(\eta_t\gamma_t^{-1}(x),t).                \tag{A.137.11}
\]
The path $\eta_t\gamma_t^{-1}$ starts at the identity and ends at
the PL map $ba^{-1}$. Equations (A.137.7)--(A.137.9) show that
\[
                 \kappa(h_\eta)-\kappa(h_\gamma)\in J_B.
                                                               \tag{A.137.12}
\]
Conversely every element of $J_B$ can be added by postcomposing
$h_\gamma$ with a trace (A.137.8), still in the same free-top
component and still with PL top.

The resulting coset depends only on the component of $h$, not on
the representative. To check this explicitly, it suffices after
the endpoint corrections to consider a path $H_s$ with endpoints
$F_0,F_1\in D_B$. Let $a=e(F_0)$ and $f_s=e(H_s)$, and put
\[
 T_s(x,t)=(f_{s t}a^{-1}(x),t),\qquad Q_s=T_s^{-1}H_s.
                                                               \tag{A.137.13}
\]
Every $Q_s$ has the same boundary maps, identity at the bottom and
$a$ at the top. Apply the common-collar normalization to this
family; its fixed PL boundary maps make this a family of relative
PL structures. Thus $\kappa(Q_s)$ is constant. At $s=0$ it is
$\kappa(F_0)$, while $T_1$ is the trace of a path from the identity
to the PL map $e(F_1)a^{-1}$. At $s=1$, additivity gives
\[
        \kappa(F_1)=\kappa(T_1)+\kappa(F_0).
\]
Its difference lies in $J_B$. This proves that $\lambda_B$ in the
statement is well-defined. The formula also verifies the assertion
for arbitrary continuous paths of homeomorphisms, rather than only
block pseudo-isotopies.

The set $I_B$ is a subgroup, being the inverse image of a subgroup
under $\pi_0e$. Given two classes in $I_B$, choose representatives
in $D_B$ by (A.137.10). Their product represents the product of the
original classes and still belongs to $D_B$. Additivity (A.137.7)
proves that $\lambda_B$ is a homomorphism.

\textbf{The exact image on components.}
A PL concordance has PL top and $\kappa=0$, so its component lies
in $\ker\lambda_B$. Conversely suppose $[h]\in I_B$ and
$\lambda_B([h])=0$. Choose $h_\gamma\in D_B$ as above.
Its class $\kappa(h_\gamma)$ belongs to $J_B$. Choose a trace
$T_\delta$ whose class is its additive inverse, using (A.137.9).
Then $T_\delta h_\gamma$ lies in the same free-top component as
$h$ and has $\kappa=0$ by (A.137.7). The zero test (A.137.4)
isotopes it, with both boundary maps now fixed, to a PL concordance.
This proves (A.137.2). A component outside $I_B$ cannot contain a
PL concordance because top restriction preserves paths. If $H^2(B)$
vanishes the quotient obstruction group is zero, giving the stated
special case. Empty products give the empty-base assertion.
The first isomorphism theorem gives (A.137.2a). A compact finite PL
manifold has finite-dimensional mod-$2$ cohomology, so its target,
and hence the displayed quotient, is a finite elementary abelian
$2$-group with the asserted bound.

The proof used relative PL structures in ambient dimension five,
always after both endpoint maps had become PL. It never replaced
the first endpoint condition by stable PL data on $B$, nor used a
smooth-to-PL comparison in dimension four. The explicit ordinary
isotopies retain the free top precisely where it is permitted.
\end{proof}
